%% Chapter 1 — Introduction
\chapter{Introduction}
\label{ch:introduction}

\section{Motivation}
\label{sec:motivation}

Communication through avatars has become a defining feature of virtual reality, extending well beyond its origins as a navigational convenience to serve as the primary medium through which users express identity, signal intent, and form social relationships in shared digital environments.
The perceptual consequences of this shift are well documented.
Avatar appearance exerts a measurable influence on social presence and interpersonal trust, meaning that the visual design of an avatar functions as a communicative act in its own right, shaping how users are perceived and how they engage with one another~\cite{nowak2003, weidner2023}.

Recent advances in generative AI have made photorealistic avatar creation more accessible than ever.
Platforms such as Avaturn can now produce a fully rigged, expressively animated CG avatar from a single uploaded photograph, reconstructing facial geometry and surface appearance with a degree of fidelity that was difficult to achieve outside of dedicated production pipelines only a few years ago.
The implicit promise of this technology is that a digital representation close enough to the real person will be received as one.
Yet this is precisely where a well-known perceptual problem arises.
Mori's uncanny valley hypothesis, first proposed in 1970, predicts that as a human representation approaches but falls just short of true photorealistic fidelity, it tends to evoke unease rather than familiarity~\cite{mori1970}.
An avatar that comes close to realism without quite reaching it can therefore undermine the very trust it was designed to support, a finding corroborated across robotics, computer graphics, and virtual environment research.

Non-photorealistic rendering~(NPR) offers a different path through this problem.
Rather than simulating accurate lighting and surface detail, NPR techniques deliberately stylise the rendered output to achieve artistic or communicative goals~\cite{gooch1998}.
Cel shading, edge line drawing, and painterly filtering can each abstract an avatar's appearance in ways that sidestep the uncanny valley, while preserving enough expressiveness and readable identity that the character remains engaging.
The open question is not whether this looks appealing.
It is whether a deliberately abstracted avatar can still be trusted by the people interacting with it, or whether reducing realism introduces its own perceptual costs.
That question has received surprisingly little attention in the context of modern, SDK-managed avatar systems.

\section{Problem Statement}
\label{sec:problem_statement}

Applying NPR stylisation to photorealistic avatar systems is not straightforward.
Production avatar platforms each manage rendering through their own shader architectures, material systems, and mesh topologies, and integrating custom visual effects into these pipelines without disrupting avatar fidelity or performance requires working within constraints that most prior NPR research has not had to address.
Existing work on NPR for virtual characters largely operates on custom assets built under full authorial control.
How NPR techniques behave when applied across platforms with different rendering architectures, and which constraints prove most limiting in practice, is not well established.

Beyond this technical question, there is an open empirical one.
Photorealistic avatar generation has become increasingly accessible, yet the uncanny valley problem suggests that realism alone does not guarantee a positive reception.
Whether NPR abstraction can improve how an avatar is perceived, specifically in terms of trustworthiness, and how much stylisation is too much before credibility is lost, are questions that remain largely unanswered.
They matter for any application where avatar credibility is at stake: social VR, remote consultation, education, and more.

This thesis addresses both challenges.
It examines how NPR techniques can be applied across avatar platforms with differing rendering constraints, and asks whether the resulting visual abstraction affects the trust that viewers place in an avatar delivering a recommendation.

\section{Contributions}
\label{sec:contributions}

This thesis makes the following contributions:

\begin{itemize}
  \item \textbf{Multi-platform NPR shader system.}
    A real-time NPR stylisation system is developed and applied across three avatar platforms: a Mixamo humanoid model used for initial technique development, the Meta Avatars SDK integrated via the app-specific post-lighting fragment hook~(\texttt{AppSpecificPostManipulation}), and Avaturn, which serves as the primary platform for the user study.
    The cross-platform implementation reveals the constraints and trade-offs specific to each rendering pipeline and demonstrates that the core NPR techniques generalise across different mesh topologies and material systems.

  \item \textbf{Iteratively developed stylisation techniques.}
    A series of techniques are designed, implemented, and compared through an iterative development process:
    toon shading with a quantized diffuse model,
    X-Toon extended toon shading with a depth-indexed 2D ramp,
    screen-space normal edge detection,
    Sobel edge detection with adaptive skin-to-clothing thresholding,
    Gaussian pre-filtered Sobel,
    hierarchical multi-cue edge detection combining depth, normal, and colour discontinuities,
    isotropic Kuwahara painterly filtering,
    Kuwahara combined with Sobel edge lines,
    and a composite technique combining Kuwahara, Gaussian Sobel, and hierarchical detection.
    Each technique addresses limitations observed in its predecessor, and all are selectable at runtime without recompilation.
    A final technique is selected for the user study based on visual assessment across all three avatar platforms.

  \item \textbf{In-VR parameter tuning interface.}
    A fully procedural, WorldSpace UI panel is implemented in Unity, enabling live exploration of all shader parameters from inside the headset using controller raycasting.
    This tool was central to the iterative development process and can be reused for future shader experimentation.

  \item \textbf{Video-based user study on avatar abstraction and trust.}
    A user study is conducted using pre-recorded video scenarios in which a speaker, shown as either a real person, a photorealistic Avaturn avatar, or an NPR-stylised avatar, recommends one of two options in a practical decision situation.
    The video format is chosen deliberately: it is the only medium that allows a real human speaker and avatar counterparts to be compared under identical viewing conditions.
    The study provides empirical insight into how visual abstraction affects the perceived trustworthiness of an avatar's advice.
\end{itemize}

\section{Thesis Outline}
\label{sec:outline}

The remainder of this thesis is structured as follows:

\begin{description}
  \item[\Cref{ch:related_work}] surveys prior work on NPR for 3D characters, avatar rendering and realism in VR, and trust and perception in virtual environments, positioning this thesis in the existing research landscape.

  \item[\Cref{ch:background}] introduces the technical and theoretical foundations required for this work, covering edge detection mathematics, the Kuwahara filter, the Meta Avatars SDK rendering pipeline, and concepts of trust and abstraction from human-computer interaction.

  \item[\Cref{ch:methodology}] presents the research questions, the overall system design, the iterative development strategy, and the user study design.

  \item[\Cref{ch:implementation}] describes the full implementation of the NPR shader system, covering all eight techniques in detail, the runtime tooling, and the rationale for the final technique selection used in the user study.

  \item[\Cref{ch:evaluation}] reports the technical evaluation of each stylisation technique and presents the results of the user study, followed by a discussion of the findings and their implications.

  \item[\Cref{ch:conclusion}] summarises the contributions of this thesis, reflects on its limitations, and outlines directions for future work.
\end{description}
